%! Confidence Intervals for the Difference of Two Means — Unit 7, Lesson 7.6 — Group Activity (Answer Key)
\documentclass[10pt]{article}
\usepackage{apstats-article}
\usepackage{apstats-key}   % pulls in -boxes; do NOT also load -boxes
\newcommand{\TallMath}[1]{$\displaystyle #1\rule[-1.4em]{0pt}{3.2em}$}

\begin{document}

\pageheader{Unit 7, Lesson 7.6}{Group Activity --- Answer Key}
\namepartnerperiod

\vspace{0.15in}

\textit{For each scenario: verify conditions, compute the standard error, find the confidence
interval, and interpret the result. Show all work.}

\begin{scenariobox}[Scenario 1 --- Sleep Hours: Teens vs.\ Adults (Tier R)]{sky}
A health researcher randomly selects 32 teenagers and 35 adults and records their average
nightly sleep (hours). Results:
Teens: $\bar x_1 = 7.1$ hr, $s_1 = 1.2$ hr. \quad
Adults: $\bar x_2 = 7.8$ hr, $s_2 = 0.9$ hr.
Both sample distributions are roughly symmetric with no outliers.

\begin{enumerate}[label=\alph*.]
  \item \textbf{Identify the procedure} and write the point estimate.

        Procedure: \ans{two-sample $t$-interval for $\mu_{\text{teens}} - \mu_{\text{adults}}$}

        Point estimate: $\bar x_1 - \bar x_2 = {\color{keyred}\mathbf{7.1}} - {\color{keyred}\mathbf{7.8}} =$ \ans{$-0.7$ hr}

  \item \textbf{Verify conditions.}
        \begin{itemize}
          \item Random: \ansline{Random samples of 32 teens and 35 adults --- independence condition met.}
          \item Normal: \ansline{$n_1 = 32 \ge 30$ and $n_2 = 35 \ge 30$; CLT applies --- Normal condition met for both.}
        \end{itemize}

  \item \textbf{Compute} the standard error.

        $SE = \sqrt{\dfrac{s_1^2}{n_1} + \dfrac{s_2^2}{n_2}}
            = \sqrt{\dfrac{({\color{keyred}\mathbf{1.2}})^2}{{\color{keyred}\mathbf{32}}} + \dfrac{({\color{keyred}\mathbf{0.9}})^2}{{\color{keyred}\mathbf{35}}}}
            = \sqrt{{\color{keyred}\mathbf{0.04500}} + {\color{keyred}\mathbf{0.02314}}} \approx$ \ans{0.2610}

  \item Using technology: $df \approx 61$, $t^* \approx 2.000$ (95\% confidence).

        $ME = t^* \times SE = 2.000 \times {\color{keyred}\mathbf{0.2610}} \approx$ \ans{0.522}

        $CI = ({\color{keyred}\mathbf{-0.7}} \pm {\color{keyred}\mathbf{0.522}}) = ($ \ans{$-1.222$} , \ans{$-0.178$} $)$

  \item \textbf{Interpret} the interval in context.
        \ansline{We are 95\% confident the true difference in mean nightly sleep (teens $-$ adults) is between $-1.22$ and $-0.18$ hours.}

  \item Does the interval contain 0? \ans{No} \quad
        What does this suggest? \ansline{The interval is entirely negative, providing convincing evidence that teens sleep less than adults on average ($\mu_{\text{teens}} < \mu_{\text{adults}}$).}
\end{enumerate}
\end{scenariobox}

\begin{scenariobox}[Scenario 2 --- Screen Time: Urban vs.\ Rural Students (Tier A)]{sky}
An education researcher randomly samples students in urban and rural schools to compare
mean daily screen time (hours). Results:
Urban: $n_1 = 28$, $\bar x_1 = 5.4$ hr, $s_1 = 1.6$ hr. \quad
Rural: $n_2 = 25$, $\bar x_2 = 4.1$ hr, $s_2 = 1.3$ hr.
Both sample distributions are roughly symmetric with no outliers.

\begin{enumerate}[label=\alph*.]
  \item \textbf{Verify both conditions.}
        \begin{itemize}
          \item Random / Independence: \ansline{Random samples; two independent groups --- met.}
          \item Normal: \ansline{$n_1 = 28 < 30$ and $n_2 = 25 < 30$; both sample distributions are roughly symmetric with no outliers --- Normal condition met for each.}
        \end{itemize}

  \item \textbf{Compute} $SE$ and construct a 95\% confidence interval. $df \approx 50$, $t^* \approx 2.009$.

        $SE = \sqrt{\dfrac{({\color{keyred}\mathbf{1.6}})^2}{{\color{keyred}\mathbf{28}}} + \dfrac{({\color{keyred}\mathbf{1.3}})^2}{{\color{keyred}\mathbf{25}}}} \approx$ \ans{0.3967}

        $CI = ({\color{keyred}\mathbf{1.3}} \pm {\color{keyred}\mathbf{2.009}} \times {\color{keyred}\mathbf{0.3967}}) = ($ \ans{0.503} , \ans{2.097} $)$

  \item \textbf{Interpret} in context.
        \ansline{We are 95\% confident the true difference in mean daily screen time (urban $-$ rural) is between 0.503 and 2.097 hours.}
        \ansline{Since 0 is not in the interval and the interval is entirely positive, this provides convincing evidence that urban students spend more time on screens than rural students, on average.}
\end{enumerate}
\end{scenariobox}

\begin{scenariobox}[Scenario 3 --- Compare Two Intervals (Tier E)]{sky}
A researcher computes two separate 95\% confidence intervals:
\begin{itemize}
  \item Group comparison A: $\mu_1 - \mu_2$: $(-0.3, \; 4.1)$ minutes
  \item Group comparison B: $\mu_3 - \mu_4$: $(2.8, \; 9.4)$ minutes
\end{itemize}

\begin{enumerate}[label=\alph*.]
  \item For interval A, does the data provide convincing evidence that $\mu_1 \ne \mu_2$?
        \ansline{No. Interval A contains 0 ($-0.3$ to $4.1$), so we cannot rule out $\mu_1 = \mu_2$. There is not convincing evidence of a difference.}

  \item For interval B, does the data provide convincing evidence that $\mu_3 > \mu_4$?
        \ansline{Yes. Interval B is entirely positive ($2.8$ to $9.4$), which means 0 is not plausible. This provides convincing evidence that $\mu_3 > \mu_4$.}

  \item Give two possible reasons Interval B is wider than Interval A.
        \ansline{(1) Smaller sample sizes produce larger standard errors and wider intervals. (2) Higher variability (larger $s$) in one or both groups increases the standard error and widens the interval.}
\end{enumerate}
\end{scenariobox}

\begin{teachernote}
Scenario 3 develops meta-level thinking: students must reason about what the interval
implies rather than just computing it. The two reasons for interval width connect back to
the SE formula $\sqrt{s_1^2/n_1 + s_2^2/n_2}$ --- bigger $s$ or smaller $n$ both increase $SE$.
\end{teachernote}

\end{document}
